\documentclass[11pt,a4paper]{article}
\usepackage[utf8]{inputenc}
\usepackage[T1]{fontenc}
\usepackage[french]{babel}
\usepackage[margin=2.5cm]{geometry}
\usepackage{listings}
\usepackage{xcolor}
\usepackage{hyperref}

\lstset{
  basicstyle=\ttfamily\small,
  backgroundcolor=\color{gray!8},
  frame=single,
  breaklines=true,
  columns=fullflexible,
  showstringspaces=false
}

\title{TP3 de compilation\\Initiation à Flex, générateur d'analyseur lexical}
\author{Nom Prénom -- ILISI 2}
\date{\today}

\begin{document}
\maketitle

\section{Principe de Flex}

Un fichier Flex (\texttt{.l}) a trois parties séparées par \texttt{\%\%} :
\begin{lstlisting}
definitions (options, macros, conditions de depart, code C %{ ... %})
%%
regles :   expression_reguliere    { action en C }
%%
code C (fonction main, fonctions utilitaires)
\end{lstlisting}

Flex transforme les expressions régulières en un automate fini (fichier \texttt{lex.yy.c}).
Le fonctionnement de \texttt{yylex()} repose sur deux règles :
\begin{itemize}
  \item \textbf{plus longue correspondance} : on choisit la règle qui reconnaît le texte le plus long ;
  \item \textbf{priorité} : à longueur égale, la règle écrite en premier gagne.
\end{itemize}
Le texte reconnu est dans \texttt{yytext}. Tout caractère non reconnu est recopié
tel quel (règle par défaut \texttt{ECHO}).

\paragraph{Compilation.}
\begin{lstlisting}[language=bash]
flex exercice.l
gcc lex.yy.c -o exercice
./exercice < fichier.txt
\end{lstlisting}

\section{Exercice 1.1 : programme source d'un langage simple}

\subsection{a) Reconnaissance de symboles}

\begin{lstlisting}
[a-zA-Z][a-zA-Z0-9]*    { printf("<%s> ", yytext); }
\end{lstlisting}

\textbf{Logique.} Un symbole commence par une lettre, suivie de zéro ou plusieurs
lettres ou chiffres. Pour \texttt{for( i=0; i < 10; i++ )}, on obtient
\texttt{<for>( <i>=0; <i> < 10; <i>++ )}.
Les nombres (\texttt{10}, \texttt{12.3}, \texttt{.0}) ne commencent pas par une lettre ;
les ponctuations et espaces non plus. Ils sont donc recopiés sans changement par la règle par défaut.
Le mot-clé \texttt{for} est traité comme un symbole ordinaire.

\subsection{b1) Mots-clés, opérateurs, symboles}

\begin{lstlisting}
for|while|if|else|do|int|float|...  { printf("<MOT-CLE:%s> ", yytext); }
[a-zA-Z][a-zA-Z0-9]*                { printf("<SYMBOLE:%s> ", yytext); }
"++"|"--"|"=="|"<="|">="|"&&"|...   { printf("<OPERATEUR:%s> ", yytext); }
[-+*/%=<>!&|^~]                     { printf("<OPERATEUR:%s> ", yytext); }
[(){}\[\];,.:]                      { printf("<OPERATEUR:%s> ", yytext); }
\end{lstlisting}

\textbf{Logique.} Pour le mot \texttt{for}, la règle des mots-clés et celle des symboles
reconnaissent la même longueur. Par priorité, la règle des mots-clés doit donc être
\emph{écrite avant}. Pour \texttt{++}, la règle à deux caractères est plus longue que deux
fois la règle \texttt{+} : \texttt{++} est reconnu comme un seul opérateur.

\subsection{b2) Valeurs numériques}

\begin{lstlisting}
EXP     [eE][-+]?[0-9]+
%%
[0-9]+                                                  { printf("<ENTIER:%s> ", yytext); }
[0-9]+"."[0-9]*({EXP})?|"."[0-9]+({EXP})?|[0-9]+{EXP}   { printf("<FLOTTANT:%s> ", yytext); }
\end{lstlisting}

\textbf{Logique.} Un entier est une suite de chiffres. Un flottant a une partie
décimale (\texttt{12.3}, \texttt{12.}, \texttt{.0}) et éventuellement un exposant
(\texttt{1e5}). Pour \texttt{12.3}, la règle flottant (4 caractères) l'emporte sur
l'entier \texttt{12} (2 caractères) : c'est la plus longue correspondance.

\section{Exercice 1.2 : suppression des commentaires C}

\begin{lstlisting}
%x COMMENT
%{ int nb_comm = 0, nb_lignes = 0; %}
%%
\"([^"\\\n]|\\.)*\"    { ECHO; }
"/*"                   { BEGIN(COMMENT); nb_comm++; }
<COMMENT>"*/"          { BEGIN(INITIAL); }
<COMMENT>\n            { nb_lignes++; }
<COMMENT>.             { /* ignore */ }
<COMMENT><<EOF>>       { fprintf(stderr, "commentaire non termine\n"); return 0; }
\end{lstlisting}

\textbf{Logique.} Un commentaire peut s'étendre sur plusieurs lignes : une seule
expression régulière est possible, mais on utilise ici une \emph{condition de départ}
(\texttt{\%x COMMENT}), plus lisible.
\begin{itemize}
  \item En état normal, le texte est recopié. À la lecture de \texttt{/*}, on passe dans l'état \texttt{COMMENT}.
  \item Dans cet état, tout est ignoré (aucune sortie), sauf \texttt{*/} qui ramène à l'état normal.
  \item \textbf{Variables dans les actions} : \texttt{nb\_comm} et \texttt{nb\_lignes} sont des
        variables C globales modifiées dans les actions et affichées à la fin dans \texttt{main}.
  \item La règle des chaînes \texttt{"..."} protège un \texttt{/*} qui serait à l'intérieur d'un texte.
  \item \texttt{<<EOF>>} dans l'état \texttt{COMMENT} signale un commentaire jamais fermé.
\end{itemize}

\section{Exercice 1.3 : expressions régulières}

Hypothèse : un mot par ligne. Les délimiteurs (espace, tabulation, fin de ligne) sont
simplement recopiés ; une règle « attrape-tout » (\texttt{.}) signale chaque caractère
d'un mot rejeté.
Chaque règle est \emph{ancrée} avec \texttt{\string^} (début de ligne) et \texttt{\$} (fin de ligne),
pour que le mot entier soit vérifié et non un morceau de mot.

Structure commune :
\begin{lstlisting}
\n                  { ECHO; }
[ \t]               { ECHO; }
^ ER $              { printf("%s  -> ACCEPTE", yytext); }
.                   { printf("<?%c>", yytext[0]); }
\end{lstlisting}

\subsection{a) Les 5 voyelles dans l'ordre}
\begin{lstlisting}
^{L}*[aA]{L}*[eE]{L}*[iI]{L}*[oO]{L}*[uU]{L}*$       (L = [a-zA-Z])
\end{lstlisting}
\textbf{Logique.} On autorise n'importe quelles lettres entre les voyelles :
\texttt{a}, puis \texttt{e}, puis \texttt{i}, \texttt{o}, \texttt{u}, dans cet ordre.
Exemples acceptés : \texttt{abstemious}, \texttt{aeiou}. Rejeté : \texttt{sequoia}.

\subsection{b) Commence ou finit par \texttt{a}}
\begin{lstlisting}
^(a[ab]*|[ab]*a)$
\end{lstlisting}
\textbf{Logique.} Deux cas reliés par \texttt{|} : \texttt{a} suivi de n'importe quoi,
ou n'importe quoi suivi de \texttt{a}.

\subsection{c) Pas deux caractères identiques consécutifs}
\begin{lstlisting}
^((ab)+a?|(ba)+b?|a|b)$
\end{lstlisting}
\textbf{Logique.} Si deux lettres consécutives sont toujours différentes, le mot
\emph{alterne} : \texttt{abab...} ou \texttt{baba...}. On écrit donc ces deux familles
(avec une lettre finale optionnelle), plus les mots d'une lettre.

\subsection{d) Ne contient pas la sous-chaîne \texttt{aba}}
\begin{lstlisting}
^(b+)?(a+bb+)*(a+b?)?$
\end{lstlisting}
\textbf{Logique.} \texttt{aba} signifie qu'un seul \texttt{b} est enfermé entre deux
\texttt{a}. Il faut donc que tout bloc de \texttt{b} situé entre deux blocs de \texttt{a}
contienne au moins deux \texttt{b}. Le mot se décompose en :
\begin{enumerate}
  \item des \texttt{b} au début (\texttt{(b+)?}) ;
  \item des groupes « bloc de a puis au moins deux b » (\texttt{(a+bb+)*}) ;
  \item un dernier bloc de \texttt{a}, éventuellement suivi d'un seul \texttt{b} en fin de mot (\texttt{(a+b?)?}).
\end{enumerate}
Accepté : \texttt{abba}, \texttt{bbaa}. Rejeté : \texttt{aba}, \texttt{abab}, \texttt{baba}.

\section{Programme avec menu}

Le fichier \texttt{tp3\_menu.l} regroupe tous les exercices dans un seul analyseur.
Chaque exercice correspond à une condition de départ (\texttt{\%x E11A E11B1 E11B2 E12 E13A\dots}),
et les règles sont préfixées par la condition, par exemple \texttt{<E13B>\string^(a[ab]*|[ab]*a)\$}.
Dans \texttt{main}, un menu lit le choix de l'utilisateur, charge le texte (fichier ou saisie clavier),
l'envoie à l'analyseur avec \texttt{yy\_scan\_string}, active la condition par \texttt{BEGIN(...)}
puis appelle \texttt{yylex()}.

\section{Conclusion}

Ce TP montre qu'avec Flex, une analyse lexicale se réduit à écrire des expressions
régulières associées à des actions C. Les trois notions clés sont : la plus longue
correspondance, la priorité de l'ordre des règles, et les conditions de départ
pour les analyses dépendant du contexte (commentaires).

\end{document}
